

\section{Les espaces de nom (pénible, mais parfois utile)}


\begin{frame}{Espaces de nom (namespaces)}

Les espaces de noms permettent de :
\begin{itemize}
\item Lever les ambiguités sur les noms d'éléments et d'attributs
\begin{itemize}
\item Si \verb?nom\_elt? de $doc_1$ et \verb?nom\_elt? de $doc_2$ ne représentent pas la m\^eme information
\item Si \verb?nom\_elt1? de $doc_1$ et \verb?nom\_elt2? de $doc_2$ représentent pas la m\^eme information mais sont définis différemment.
\end{itemize}
\item Fonction de regroupement des éléments
\end{itemize}

$\rightarrow$ De fa\c{c}on plus générale, sert à interpréter correctement les éléments. 

\prof{cas1 : pe parachute (classique ou doré ?)
cas 2 : pe adresses (\`a plat) ou découpées en sous-éléments
Tr\`es utile pour l'intégration de docs issus du web (domaines différents, vocabulaires différents)
Permet également une tra\c{c}abilité des balises
M\^eme principe que les packages
}
\end{frame}


% \begin{frame}{Espaces de nom (namespaces)}

% \begin{tikzpicture}[%
%   back line/.style={densely dotted},
%   cross line/.style={preaction={draw=white, -,line width=6pt}}]
%   \node (A) {};
%   \node [right of=A] (B) {};
%   \node [below of=A] (C) {};
%   \node [right of=C] (D) {};
 
%   \node (A1) [right of=A, above of=A, node distance=0.5cm] {};
%   \node [right of=A1] (B1) {};
%   \node [below of=A1] (C1) {};
%   \node [right of=C1] (D1) {};
 
% %  \draw[back line] (D1) -- (C1) -- (A1);
%  % \draw[back line] (C) -- (C1);
%   \draw[cross line] (D1) -- (B1) -- (A1) -- (A)  -- (B) -- (D) -- (C) -- (A);
%   \draw (D) -- (D1) ;
%   \draw (D1) -- (B1) -- (B) ;


%   \node (A2) [right of=A, node distance=5cm] {toto};

% \end{tikzpicture}

% \end{frame}


 
\begin{frame}{Espaces de nom (\textit{namespaces})}

\begin{center}
\figSized{images/ns1.jpg}{Espace de noms pour associer une sémantique aux éléments}{0.9}
\end{center}

\end{frame}


\begin{frame}{Espaces de nom (namespaces)}

\begin{center}
\figSized{images/ns2.jpg}{Espace de noms pour associer une sémantique aux éléments}{0.9}
\end{center}

\end{frame}


\begin{frame}{Espaces de nom (namespaces)}

\begin{center}
\figSized{images/ns3.jpg}{Espace de nom pour associer une sémantique aux éléments}{0.9}
\end{center}

\end{frame}


\begin{frame}{Espaces de nom (namespaces)}

\begin{center}
\figSized{images/ns4.jpg}{Espace de nom pour associer une information de provenance aux éléments}{0.8}
\end{center}


\end{frame}



\begin{frame}{Déclaration}

Il existe deux façons d'associer un espace de noms à un (ou plusieurs éléments) dans un fichier XML :
\begin{itemize}
\item par défaut, ou
\item explicitement.
\end{itemize}

\end{frame}

\begin{frame}[containsverbatim]{Déclaration par défaut}

\begin{lstlisting}
<genie xmlns=''http://chezmoi.fr/''>
   <prenom>Alan</prenom>
   <nom>Turing</nom>
</genie>
\end{lstlisting}

\verb?http://chezmoi.fr/? est le nom (identifiant) de l'espace de noms.\\

\smallskip

L'\textbf{élément genie} et \textbf{tous ses fils} sont dans l'espace de nom \url{http://chezmoi.fr/}.

\vfill

Les textes (Alan et Turing) ne sont dans aucun espace de nom puisque ce sont des \emph{données}.
Ne s'applique pas non plus aux noms d'attributs et valeurs d'attributs.

\vfill

URI (Uniform Resource Identifier) : ce sont généralement des URL. Intérêt\,: si vous utilisez un nom de domaine qui vous appartient, garantie d'unicité.\\
\prof{URL (Locator) garanti sans ambiguté par des agences chargées de distribuer les noms de domaine par niveau (par l'AFNIC -Association française pour le nommage Internet en coopération- pour exemple .fr et .re). Utilise le protocol http Pour la récupération de données sur internet. URL bien adapté car un propriétaire est associé à une URL. 
	http://www.w3.org/1999/xhtml
URN (Name) garanti sans ambiguite par des normes orientées métier par exemple, numéro ISBN garanti par l'AFNIL (Agence Francophone pour la Numérotation Internationale du Livre)
urn:ISBN:0-395-36341-6.\\
L'important est que ce soit un identifiant. D'autres protocoles peuvent être utilisés :
mailto:virginie.thion@cnam.fr
urn:isbn:0-201-53771-0
geo:27.988056;86.925278}

\vfill

Cf. packages Java.

\prof{URL bien adapté car un propriétaire est associé à une URL. 
Petite subtilité : l'URL ne signifie pas qu'un document existe \`a l'adresse indiquée. Ce n'est qu'un identifiant.
}

\end{frame}

\begin{frame}[containsverbatim]{Traduction par un interpréteur}

\begin{lstlisting}
<genie xmlns="http://chezmoi.fr/">
   <prenom> Alan </prenom>
   <nom> Turing </nom>
</genie>
\end{lstlisting}

Ce que comprend (traduit) un interpréteur :

\begin{lstlisting}
<http://chezmoi.fr/genie>
   <http://chezmoi.fr/prenom> Alan </http://chezmoi.fr/prenom>
   <http://chezmoi.fr/nom> Turing </http://chezmoi.fr/nom>
</http://chezmoi.fr/genie>
\end{lstlisting}

\bigskip

Le \emph{nom étendu} \verb?http://chezmoi.fr/genie? est composé 
\begin{itemize}
\item du \emph{nom de l'espace de noms} associé à l'élément (\verb?http://chezmoi.fr/?) et 
\item du \emph{nom local} de l'élément (\verb?genie?).
\end{itemize}

\end{frame}


\begin{frame}[containsverbatim]{Example de déclaration par défaut}

\begin{lstlisting}[caption={Utilisation de l'espace de nom de MathML}]
<msup xmlns="http://www.w3.org/1998/Math/MathML">
    <mfenced>
        <mrow>
          <mi>a</mi>
          <mo>+</mo>
          <mi>b</mi>
        </mrow>
      </mfenced>
      <mn>2</mn>
    </msup>
\end{lstlisting}

\begin{exoblock}
Comment ce code est-il traduit par un interpréteur ?
\end{exoblock}

\prof{ \tiny
\verb?<http://www.w3.org/1998/Math/MathML/msup>?\\
\verb?  <http://www.w3.org/1998/Math/MathML/mfenced>?\\
\verb?     <http://www.w3.org/1998/Math/MathML/mrow>?\\
\verb?        <http://www.w3.org/1998/Math/MathML/mi>a</http://www.w3.org/1998/Math/MathML/mi>?\\
\verb?        <http://www.w3.org/1998/Math/MathML/mo>+</http://www.w3.org/1998/Math/MathML/mo>?\\
\verb?        <http://www.w3.org/1998/Math/MathML/mi>b</http://www.w3.org/1998/Math/MathML/mi>?\\
\verb?     </http://www.w3.org/1998/Math/MathML/mrow>?\\
\verb?  </http://www.w3.org/1998/Math/MathML/mfenced>?\\
\verb?  <http://www.w3.org/1998/Math/MathML/mn>2</http://www.w3.org/1998/Math/MathML/mn>?\\
\verb?</http://www.w3.org/1998/Math/MathML/msup>?\\
}

\prof{Petite subtilité : l'URL ne signifie pas qu'un document existe à l'adresse indiquée. Ce n'est qu'un identifiant.
}
\end{frame}

\begin{frame}[containsverbatim]{Déclaration explicite}

\begin{lstlisting}[caption=Déclaration explicite d'un espace de nom]
<genie xmlns:p=''http://chezmoi.fr/''>
   <prenom>Alan</prenom>
   <nom>Turing</nom>
</genie>
\end{lstlisting}
Dans le code ci-dessus, l'espace de nom est \alert{déclaré} mais n'est \alert{appliqué} à aucun élément ni attribut. \\

\vfill

La déclaration comprend un \alert{préfixe}, ``\alert{p}'': une fois déclaré, il est employable dans tout élément. 

\vfill
Si un élément ou un attribut n'est pas prefixé, il n'est pas dans l'espace de nom.\\

\begin{lstlisting}[caption=Déclaration explicite et application d'un espace de nom]
<p:genie xmlns:p=http://chezmoi.fr/ p:date_naiss=''p:1912-06-23''>
   <p:prenom>Alan</p:prenom>
   <nom>Turing</nom>
</p:genie>
\end{lstlisting}

\prof{Montrer : application à un élément, à un attribut, à une valeur d'attribut, et non application\\
Explicite pour \\
1) plus de souplesse (car les espaces de noms car par défaut, s'applique à tous les sous-éléments, ce qu'on ne désire pas forcément. )\\
2) pouvoir appliquer les espaces de noms aux attributs et valeurs d'attributs
}

\end{frame}


\begin{frame}[containsverbatim]{Déclaration explicite}

\begin{lstlisting}
<p:genie xmlns:p=http://chezmoi.fr/ p:date_naiss="p:1912-06-23">
   <p:prenom> Alan </p:prenom>
   <nom> Turing </nom>
</p:genie>
\end{lstlisting}

Ce que traduit l'interpréteur :

\begin{lstlisting}
<http://chezmoi.fr/genie http://chezmoi.fr/date_naiss="http://chezmoi.fr/1912-06-23">
   <http://chezmoi.fr/prenom> Alan </http://chezmoi.fr/prenom>
   <nom> Turing </nom>
</http://chezmoi.fr/genie>

\end{lstlisting}
\end{frame}

